\section{引言}\label{sec:intro}

对赝标介子结构的研究，有助于揭示量子色动力学（QCD）中由来已久的质量起源之谜。
在标准模型中，希格斯玻色子提供了所有非复合费米子质量的绝大部分。
然而，介子价夸克的质量仅贡献介子总质量的一小部分。
其余部分被认为源自强相互作用的动力学。
研究 K 介子与 $\pi$ 介子的结构，能极大地帮助我们理解强子质量的涌现。
一类常被研究的结构是部分子分布函数（PDF），它给出在强子内部找到夸克与胶子的概率，并表示为它们占母强子动量份额 $x$ 的函数。
由于 $\pi$ 介子的质量几乎全部来自胶子场，研究介子的胶子 PDF 对理解介子结构非常有价值。
人们已多次尝试提取 $\pi$ 介子胶子 PDF~\cite{Barry:2018ort,Cao:2021aci,Novikov:2020snp,Chang:2020rdy}，但在 K 介子 PDF 方面进展甚微。
未来的实验研究，例如美国的电子离子对撞机（EIC）~\cite{Accardi:2012qut,AbdulKhalek:2021gbh}、中国的电子离子对撞机（EicC）~\cite{Anderle:2021wcy}，
以及欧洲的 Drell--Yan 与 $J/\psi$ 产生实验 COMPASS++/AMBER~\cite{Denisov:2018unj}，都将致力于改进我们对 $\pi$ 介子与 K 介子胶子和夸克 PDF 的认识。
关于当前及未来 $\pi$ 介子与 K 介子结构研究的更多细节，读者可参阅近期综述~\cite{Aguilar:2019teb,Roberts:2021nhw,Arrington:2021biu}。


格点 QCD 为确定标准模型输入提供了一种方法，可增进我们对非微扰介子胶子结构的认识。
然而，由于胶子观测量存在众所周知的噪声—信号比问题，确定 $\pi$ 介子胶子 PDF 第一矩的工作仅有寥寥数例~\cite{Meyer:2007tm,Shanahan:2018pib}，且所计算的 $\pi$ 介子质量最轻为 450~MeV。
尚无任何 K 介子矩的计算。
我们几乎没有可能通过格点计算越来越高的矩来重构介子 PDF 的 $x$ 依赖性。
最近，人们提出了强子胶子 PDF 的直接格点计算方法：赝部分子分布（pseudo-PDF）方法~\cite{Balitsky:2019krf} 与准部分子分布（quasi-PDF）方法~\cite{Zhang:2018diq,Wang:2019tgg}，二者均源于大动量有效理论（LaMET）的提出~\cite{Ji:2013dva,Ji:2014gla,Ji:2017rah}。
采用其他方法的持续努力，例如强子张量~\cite{Liu:1993cv}、算符乘积展开~\cite{Detmold:2005gg}、格点截面~\cite{Ma:2014jla,Ma:2017pxb} 与赝 PDF 方法~\cite{Radyushkin:2017cyf,Orginos:2017kos}，在提取 $x$ 依赖的强子结构方面也取得了一些成功。
不过，大多数此类 PDF 计算采用流行的准 PDF 与赝 PDF 技术，且大多局限于核子中的同位旋矢量夸克分布以及 $\pi$ 介子、K 介子中的价夸克分布~\cite{Lin:2013yra,Lin:2014zya,Chen:2016utp,Lin:2017ani,Alexandrou:2015rja,Alexandrou:2016jqi,Alexandrou:2017huk,Chen:2017mzz,Alexandrou:2018pbm,Chen:2018xof,Chen:2018fwa,Alexandrou:2018eet,Lin:2018qky,Fan:2018dxu,Liu:2018hxv,Wang:2019tgg,Lin:2019ocg,Chen:2019lcm,Lin:2020reh,Chai:2020nxw,Bhattacharya:2020cen,Lin:2020ssv,Zhang:2020dkn,Li:2020xml,Fan:2020nzz,Gao:2020ito,Lin:2020fsj,Lin:2021brq,Zhang:2020rsx,Alexandrou:2020qtt,Alexandrou:2020zbe,Lin:2020rxa,Gao:2021hxl,Lin:2020rut,Orginos:2017kos,Karpie:2017bzm,Karpie:2018zaz,Karpie:2019eiq,Joo:2019jct,Joo:2019bzr,Radyushkin:2018cvn,Zhang:2018ggy,Izubuchi:2018srq,Joo:2020spy,Bhat:2020ktg,Fan:2020cpa,Sufian:2020wcv,Karthik:2021qwz,Chen:2018fwa,Sufian:2019bol,Izubuchi:2019lyk,Joo:2019bzr,Sufian:2020vzb,Shugert:2020tgq,Gao:2020ito}。

首个将准 PDF 方法应用于胶子 PDF 的探索性研究~\cite{Fan:2018dxu} 在含 2+1 味畴壁费米子的规范系综上采用重叠（overlap）价费米子，时空体积为 $24^3\times 64$，$a=0.1105(3)$~fm，$M_\pi^\text{sea}=330$~MeV。
胶子算符在全部时空格点上计算且统计量很高：对价夸克处于轻海夸克质量与奇异数质量的两组两点函数共进行了 207,872 次测量（分别对应 $\pi$ 介子质量 340 与 678~MeV）。
经傅里叶变换后，坐标空间中 $\pi$ 介子与核子胶子准 PDF 矩阵元之比与全局拟合进行了比较；
格点结果在大不确定度范围内与全局拟合一致。
随后借助赝 PDF 方法取得了进一步进展。
Fan 等人利用 MILC $N_f=2+1+1$ 格点上的 clover 费米子研究了核子胶子 PDF，采用单一格距 0.12~fm 和两个价 $\pi$ 介子质量：310 与 690~MeV~\cite{Fan:2020cpa}。
这是首个用格点 QCD 计算得到的 $x$ 依赖胶子 PDF。
Khan 等人使用蒸馏法构造的核子插值算符，在 2+1 味、$\pi$ 介子质量 358~MeV、格距 0.094~fm 的系综（ensemble）上提取了核子胶子 PDF~\cite{HadStruc:2021wmh}。
这两项工作都忽略了胶子算符与夸克单态部分的混合。
此后，Fan 继续利用具有多个格距与多个 $\pi$ 介子质量的 $N_f=2+1+1$ MILC 格点开展核子与 $\pi$ 介子的胶子 PDF 计算。
文献~\cite{Fan:2021bcr} 报道了首批由格点 QCD 计算的 $x$ 依赖 $\pi$ 介子 PDF，涉及两个格距（0.15 与 0.12~fm）以及 $M_\pi \approx 220$、310 与 690~MeV)；
该工作利用全局拟合的 PDF 夸克贡献来估计夸克混合对胶子 PDF 测定的影响。
然而，K 介子 PDF 计算仍局限于价夸克分布~\cite{Lin:2020ssv}。

在本工作中，我们利用赝 PDF 方法给出了首个 $x$ 依赖 K 介子胶子部分子分布的格点 QCD 计算。
我们采用两个格距（0.12 与 0.15~fm）和单一的 310~MeV $\pi$ 介子质量来研究 K 介子 PDF。
本文其余部分安排如下。
在第~\ref{sec:lattice-details}节中，我们讨论本计算所用的算符、格点模拟的数值设置，以及如何在格点上获得 K 介子胶子的基态矩阵元。
在第~\ref{sec:results}节中，我们讨论从格点计算中提取 K 介子胶子 PDF 的策略；
我们研究了不同步骤引入的系统不确定度，并考察了结果对格距与 $\pi$ 介子质量的依赖性。
